\section{Delete All Occurrences of a Key in a Doubly Linked List}
\label{sec:ll-delete-occurrences-dll}

\problemstatement{Given a doubly linked list and a key, delete every node
whose value equals the key, returning the (possibly new) head.}

\begin{patternbox}
The DLL's \texttt{prev} pointer makes this a single-pass sweep: at each
matching node, splice its neighbours together in $O(1)$ using both links,
handling head and tail as the only edge cases.
\end{patternbox}

\subsection*{Splice out each match in place}

Walk the list. When a node's value equals the key, connect its predecessor's
\texttt{next} to its successor and its successor's \texttt{prev} to its
predecessor. If the deleted node was the head, advance the head. Because a
DLL node knows both neighbours, each deletion is constant time with no
re-walk -- the exact advantage over a singly linked list.

\begin{ideabox}[Both Links Make Deletion Local]
Removing a node requires rewiring the links on \emph{both} sides. The
\texttt{prev} pointer supplies the left neighbour instantly, so an arbitrary
node -- including several in one pass -- is removed in $O(1)$ each.
\end{ideabox}

\subsection*{Algorithm}

\begin{enumerate}
  \item Walk from the head. For each node equal to the key: rewire
        \texttt{prev} and \texttt{next} of its neighbours; if it was the
        head, move the head forward.
  \item Continue to the next node.
  \item Return the head.
\end{enumerate}

\subsection*{Dry run}

\dryrun{DLL $1\leftrightarrow2\leftrightarrow2\leftrightarrow3$, key $2$.}

\begin{itemize}
  \item Delete both $2$s by splicing neighbours.
  \item Result $1\leftrightarrow3$.
\end{itemize}

\subsection*{C++ implementation}

\begin{lstlisting}
#include <bits/stdc++.h>
using namespace std;

struct Node {
    int val; Node* prev; Node* next;
    Node(int v) : val(v), prev(nullptr), next(nullptr) {}
};

Node* build(const vector<int>& a) {
    Node dummy(0), *t = &dummy;
    for (int x : a) { Node* n = new Node(x); t->next = n; n->prev = t; t = n; }
    Node* head = dummy.next; if (head) head->prev = nullptr;
    return head;
}

Node* deleteAll(Node* head, int key) {
    Node* cur = head;
    while (cur) {
        Node* nxt = cur->next;
        if (cur->val == key) {
            if (cur->prev) cur->prev->next = cur->next;
            else head = cur->next;                // deleting the head
            if (cur->next) cur->next->prev = cur->prev;
        }
        cur = nxt;
    }
    return head;
}

void print(Node* h) { for (; h; h = h->next) cout << h->val << " "; cout << "\n"; }

int main() {
    print(deleteAll(build({1, 2, 2, 3, 2}), 2));   // 1 3
    return 0;
}
\end{lstlisting}

\begin{complexitybox}
\textbf{Time Complexity.} $O(n)$ -- one pass. \textbf{Space Complexity.}
$O(1)$.
\end{complexitybox}

\begin{takeawaybox}
Deleting all matches in a DLL is a single-pass splice using both links,
with head/tail as the only special cases. This is the everyday DLL deletion
skill that LRU/LFU caches rely on.
\end{takeawaybox}
